\section{Designer’s Notes}
My main task in designing Tito was to translate into game terms a wealth of distinctive factors, many of them unique to the Balkan theater, that went into the conflict between the Axis and the Yugoslav guerrillas.

I started out with the concept that the conflict was a war within a war, within a war. The Axis were fighting the Allies on the major battlefronts, and what was happening in the Mediterranean and in the Soviet Union had indirect but important effects on the war against the guerrillas. But the guerrillas, while battling the Axis, were at each other’s throats. There were two rival guerrilla factions: the Communist-led Partisans headed by Marshal Tito, and the predominantly Serbian, pro-royalist Chetniks. Not only was the resolution of the war against the Axis at stake for both, but the political makeup of postwar Yugoslavia hung in the balance as well. And the guerrillas were not the only ones with internal differences. Germany, as head of the Axis bloc, saw its position in the Balkans seriously weakened by the collapse of Italy in September 1943 and the defection of Bulgaria to the Soviet side a year later.

The polyglot makeup of Yugoslavia also had to be factored in. Here is a country with six major ethnic groups, three languages, and three religions — a situation that made Yugoslavia a microcosm of the fragmented Balkans. This fragmentation and its significant influence on operations during the guerrilla war could not be ignored. For example, the regions of Serbia and Montenegro had to be depicted as the Chetnik strongholds they were, while the rest of the country had to be shown as more hospitable to the broader-based Peoples Liberation Army of Tito.

On top of this crazy-quilt political-ethnic situation, the Axis, after its lightning military conquest of Yugoslavia in April, 1941, superimposed a fragmentation of its own. Parts of the country were annexed outright by Germany, Italy, Bulgaria, and Hungary, and the rest of the land was carved into occupation zones overseen by the Germans, Italians, and Bulgarians. Into the mix were tossed the not inconsiderable pro-Axis ethnic elements in Yugoslavia, mainly Croats and Serbs, who formed national armies of their own. In true Balkans fashion, the occupation zones, which in certain cases cut across traditional ethnic boundaries, caused the Axis considerable problems in mounting a coordinated effort against the guerrillas.

In strategic terms, Yugoslavia was immensely important to the Germans. Some 50\% of Germany’s: oil, all its chromium, 60\% of its bauxite, and almost a quarter of its copper and antimony came from the Balkans, at whose heart lies Yugoslavia. The country stood astride Germany’s line of communications with its forces in Greece and Crete. When the tide of war turned against the Axis in the Mediterranean, Germany was constantly preoccupied with the fear of an Allied landing on Yugoslavia’s Adriatic coast.

So, with the guerrilla war of necessity a military sideshow, the Germans were faced with the problem of waging it as cheaply as possible, but for big stakes. Militarily, the units employed could be put to better use on the major fighting fronts. But strategically, a quiescent Yugoslavia would take some of the pressures off those fronts.

In the final analysis, what the Germans do in the game will not win the big war going on in the outside world. At best, the German Player can duplicate what the Germans historically did in Norway — surrender a relatively intact force when hostilities cease in 1945.

The guerrillas (in the game the brunt of their fighting is done by the Partisans, so mostly hereafter we’ll refer to them as such) must keep striking at the German war effort while trying to build an army capable of eventually taking the offensive against the Axis. As we shall see, the first goal is a vital precondition for achieving the second.

Operationally, the most striking factor is the mountainous nature of the'country and its overriding effects on movement and combat. To the Partisans, the mountains, which occupy 80\% of the country, were both friend and foe. The mountains offered refuge, but not sustenance. Thus, there is a compelling need for the Partisans to come out of hiding and attack to gain recruits, supplies, and self-confidence.

From the Axis standpoint, it was vital in the game to depict the general unwillingness of their forces to prosecute an anti-guerrilla war, and the devastating effects on the Partisans when the Axis occasionally got its act together and mounted large-scale anti-guerrilla operations. Then, too, there were the complications of Italy’s surrender, which threw the weight of the campaign onto German shoulders, and the defection of the Bulgarians, which suddenly confronted the Germans with a well-armed enemy on Yugoslavia’s eastern frontier.

And, speaking of well-armed enemies, there were, of course, the Soviets. Their appearance late in the game turns what had been a search-and-destroy campaign into a war with defined front lines. How heavily the Partisans rely on Soviet support determines how heavily they will be dominated by Moscow after the war. In game terms, too many Soviets helping for too long can be a mixed blessing. Dick Rustin
